% Source-derived chapter 02: Non-proper intersections
% Original source: geometric-bogomolov-conjecture-arbitrary-characteristics
\chapter{Non-proper intersections}
\label{geometric-bogomolov-conjecture-arbitrary-characteristics-chapter-02}
\source{geometric-bogomolov-conjecture-arbitrary-characteristics}

\begin{remark}[Source section]
\label{geometric-bogomolov-conjecture-arbitrary-characteristics-chapter-02-source}
\source{geometric-bogomolov-conjecture-arbitrary-characteristics}
\source{geometric-bogomolov-conjecture-arbitrary-characteristics}
This chapter reproduces the corresponding section of the retrieved
LaTeX source, including its definitions, statements, examples, and
proofs. It has not yet been translated into Lean.
\notready
\end{remark}

% The source section title was: Non-proper intersections
\label{sec nonproper}
\source{geometric-bogomolov-conjecture-arbitrary-characteristics}

Let $B$ be a smooth projective variety of dimension $d$ over an algebraically closed field $k$.
For every $i$-cycle $Z$ of $\Q$-coefficient of $B$,  denote by $[Z]$ its class in the Chow group
$\CH_i(B)_{\Q}$. For $\alpha\in \CH_i(B)_{\Q}$, write $\alpha\geq 0$ if $\alpha$ can be represented by an effective $i$-cycle.
The goal of this section is to prove the following technical result. 

\begin{pro}\label{prostrleqtot}
\source{geometric-bogomolov-conjecture-arbitrary-characteristics}
Let $g: B\to Y$ be a flat morphism between smooth projective varieties.
Let $X$ be a closed subvariety of $B$ such that $f=g|_X: X\to Y$ is surjective. 
Denote $e=\dim X-\dim Y$. 
Let $V$ be a closed subvariety of $Y$ which is not contained in 
$$Y_{e+1}=\{y\in Y|\,\, \dim f^{-1}(y)\geq e+1\}.$$ 
Let $Z_1, \dots, Z_n$ be all irreducible components of $f^{-1}(V)$ satisfying 
$f(Z_i)=V.$ Then $\dim Z_i=\dim V+e$ for $i=1,\dots,n$. 

Assume furthermore that $V\cap Y_{e+1}$ is finite and contained in $V^{\sm}$.
Then we have 
$$\sum_{i=1}^n m_i[Z_i]\leq g^*[V]\cdot [X]$$
in $\CH_{\dim V+e}(B)_\Q$. 
Here $m_i$ is the multiplicity of $Z_i$ in $f^{-1}(V).$ 
\end{pro}

Note that if $X$ is smooth, then the last inequality can be simplified as the inequality 
$$\sum_{i=1}^n m_iZ_i\leq f^*[V]$$
in $\CH_{\dim V+e}(X)_\Q$. 

The proposition will be used in \S\ref{sec final}. 
The remaining part of this section is to prove the proposition. 
Readers might skip the proof temporarily and move to the next section at the first time of reading this paper.

The idea to prove the proposition is to use complete intersections of hyperplane sections to bound the proper part of the intersection. See Proposition \ref{propropintboundhyp} for the result on complete intersection. 
In the following, we will start with a Bertini type of result to choose suitable hyperplane sections. 


\subsection{A Bertini type of result}

\begin{pro}\label{probertini}Let $Y$ be a projective variety over an algebraically closed field $k$.
\source{geometric-bogomolov-conjecture-arbitrary-characteristics}
Let $V$ be a closed subvariety of $Y$ of codimension $r\geq 1$ such that $Y$ is smooth at $\eta_V.$
Let $Z_1,\dots,Z_m$ be irreducible subvarieties of $Y$.
Assume that 
\begin{points}
\item[$\bullet$] $(\cup_{i=1}^m Z_i)\cap V$ is finite; 
\item[$\bullet$] $(\cup_{i=1}^m Z_i)\cap V\subseteq Y^{\sm}\cap V^{\sm}$.
\end{points}
Let $H$ be a Cartier divisor on $Y$ such that $\sO_Y(H)\otimes I_V$ is generated by global sections, where
 $I_V\subset \sO_Y$ is the ideal sheaf associated to $V$.
Every section $s\in H^0(\sO_Y(H)\otimes I_V)$ defines a divisor 
$H_s\in |H|$ via the inclusion $H^0(\sO_Y(H)\otimes I_V)\hookrightarrow H^0(\sO_Y(H)).$
Then for general elements 
$$(s_1,\dots, s_r)\in H^0(\sO_Y(H)\otimes I_V)^r,$$ 
the following holds:
\begin{points}
\item $H_{s_1}\cap\dots\cap H_{s_r}$ is a proper intersection in $Y$;
\item $H_{s_1}\cdots H_{s_r}=V+W$ where $W$ is an effective $(\dim Y-r)$-cycle such that 
$V\not\subseteq \Supp \,W$, and $W\cap Z_i$ is a proper intersection for every $i=1,\dots,m$.
\end{points}
\end{pro}


\proof
Set $H_i:=H_{s_i}$ for $i=1,\dots,r.$
Because $\sO_Y(H)\otimes I_V$ is generated by global sections and $s_1,\dots, s_r$ are general in $H^0(\sO_Y(H)\otimes I_V)$, 
$(H_1\setminus V)\cap\dots\cap  (H_r\setminus V)$ is a proper intersection in $Y\setminus V$, and 
$(H_1\setminus V)\cap\dots\cap  (H_r\setminus V)\cap (Z_i\setminus V)$ is proper intersection in $Y\setminus V$ for every $i=1,\dots,m$.

\medskip

As $V\subseteq H_i,$ the intersection 
$$H_1\cap \dots \cap H_r= V\cup ((H_1\setminus V)\cap\dots\cap  (H_r\setminus V))
$$
 is of pure codimension $r$ in $Y$.
So $H_1\cap \dots \cap H_r$ is a proper intersection. 
Write $H_{1}\cdots H_{r}=V+W$ where $W$ is an effective $(\dim Y-r)$-cycle.






\medskip
Pick $x_V\in Y^{\sm}(k)\cap V^{\sm}(k)$ and set $B:=(V\cap (\cup_{i=1}^m Z_i))\cup \{x_V\}$. It is a finite subset of $Y^{\sm}(k)\cap V^{\sm}(k).$
For every point $x\in B$, fix an isomorphism $\phi_x: (O(H)\otimes I_V)_x\to I_{V,x}.$  Because $x\in (V^{\sm}\cap Y^{\sm})(k)$, $I_{V,x}/ I_{V,x}I_x$ is a $k$-vector space of dimension $r.$ Because $\sO_Y(H)\otimes I_V$ is generated by global sections and $s_1, \dots, s_r$ are general in $H^0(\sO_Y(H)\otimes I_V)$, $\phi_x(s_1), \dots, \phi_x(s_r)$ is a $k$-basis of 
$I_{V,x}/ I_{V,x}I_x$. Then $\phi_x(s_1),\dots, \phi_x(s_r)$ generate $I_{V,x}.$ This implies that at every point $x\in B$, there is an open neighborhood $U_x$ of $x$ such that $H_1\cap \dots \cap H_r\cap U_x=V\cap U_x.$ 
So $V\not\subseteq \Supp \,W$, and $W\cap B=\emptyset.$
Because $V\cap (\cup_{i=1}^m Z_i)\subseteq B$ for every $i=1,\dots, m$, 
the intersection
$$W\cap Z_i=W\cap (Z_i\setminus V)=(H_1\setminus V)\cap\dots\cap  (H_r\setminus V)\cap (Z_i\setminus V)$$ is a proper intersection.
\endproof




\subsection{Proper part of an intersection}


\begin{lem}
\source{geometric-bogomolov-conjecture-arbitrary-characteristics}
\notready\label{lemcutampleeff}
\source{geometric-bogomolov-conjecture-arbitrary-characteristics}
Let $B$ be a smooth projective variety of dimension $d$ over an algebraically closed field $k$.
Let $H$ be a Cartier divisor on $B$ such that $\sO_B(H)$ is globally generated. 
Then for $\alpha\in \CH_i(B)_{\Q}$ with $\alpha\geq 0$, $\alpha\cdot [H] \geq 0.$
\end{lem}

\proof
We may write $\alpha=[Z]$ for an effective $i$-cycle $Z$. We can further assume that $Z$ is an integral subvariety of $B$. 
Because $\sO_B(H)$ is globally generated, after replacing $H$ by a general element in $|H|$, we may assume that $H\cap Z$ is a proper intersection.
Then $\alpha\cdot [H]=[Z]\cdot [H]\geq 0.$
\endproof




\medskip 

Let $B$ be a smooth projective variety of dimension $d$ over an algebraically closed field $k$.
Let $H_1,\dots,H_m$ be Cartier divisors of $B$. Let $X$ be a subvariety of $B.$
A \emph{proper component} $Z$ of $H_1\cap \dots \cap H_m\cap X$ is an irreducible component of the underling topological space of $H_1\cap \dots \cap H_m\cap X$ of dimension $\dim X-m.$ Write $m(Z, H_1\cap \dots \cap H_m\cap X)$ for the multiplicity of $Z$ in $H_1\cap \dots \cap H_m\cap X$, as defined by Serre using the tor-functor.
Define 
$$(H_1\cdots H_m\cdot X)^{\prop}:=\sum m(Z, H_1\cap \dots \cap H_m\cap X) [Z],$$
where the  sum takes over all proper components $Z$ of $H_1\cap \dots \cap H_m\cap X$.



\medskip





The following result is a generalization of \cite[Lemma 3.3]{Jia2021} with a similar proof.

\begin{pro}\label{propropintboundhyp}
\source{geometric-bogomolov-conjecture-arbitrary-characteristics}
Let $H_1, \dots, H_r$ be effective Cartier divisors on $B$ such that 
$\sO_B(H_1),\dots,\sO_B(H_r)$ are generated by global sections. 
Let $X$ be a subvariety of $B$.
Then we have 
$$(H_1\cdots H_r\cdot X)^{\prop}\leq H_1\cdots H_r\cdot X.$$
\end{pro}

\proof 
Let $V_1,\dots, V_m$ be the proper components of $H_1\cap \dots \cap H_r\cap X$ with multiplicities $m_1,\dots,m_r.$
For $i=1,\dots,m,$ set $\eta_i:=\eta_{V_i}$. Then $H_1\cap \dots \cap H_r\cap X$ has proper intersection at $\eta_1,\dots, \eta_m.$

Let $X_1,\dots, X_l$ be all irreducible components of $H_1\cap\dots\cap H_{r-1}\cap X$ passing through $\eta_1,\dots, \eta_m.$
For  $j=1,\dots,l$, the variety $X_j$ has dimension $\dim X-r+1$ and $X_j\not\subseteq H_r$
Assume that $X_j$ has multiplicity $n_j$ in $H_1\cap\dots\cap H_{r-1}\cap X$.

If $r=1,$ this lemma is trivial.  Now assume that $r\geq 2$. 
By induction,
 $$\sum_{j=1}^l n_jX_j\leq H_1\cdots H_{r-1}\cdot X.$$
The previous paragraph shows that $(\sum_{j=1}^l n_jX_j)\cap H_r$ is a proper intersection and $m_i$ is the multiplicity of $V_i$ in $(\sum_{i=1}^l n_jX_j)\cap H_r.$
By Lemma \ref{lemcutampleeff}, we have 
$$(H_1\cdots H_r\cdot X)^{\prop}=\sum_{i=1}^m m_iV_i\leq (\sum_{j=1}^l n_jX_j)\cap H_r\leq H_1\cdots H_{r-1}\cdot V\cdot H_r.$$
\endproof


\subsection{Strict transform}
Let $f:X\to Y$ be a surjective morphism of projective varieties over an algebraically closed field $k.$ 
Set $e:=\dim X-\dim Y$.

For every integer $l\geq e$,   
$$Y_l:=\{y\in Y|\,\, \dim f^{-1}(y)\geq l\}$$ 
is a closed subset of $Y$. 
We have $Y_{l+1}\subseteq Y_l$.
We note that $Y_{e}=Y$ and for $l\geq e+1,$
$\dim Y_l+l\leq \dim X-1$. In particular, for $l\geq \dim X$, $Y_l=\emptyset.$


\begin{lem}
\source{geometric-bogomolov-conjecture-arbitrary-characteristics}
\notready\label{lempullbackproper}
\source{geometric-bogomolov-conjecture-arbitrary-characteristics}
Let $W$ be a subvariety of $Y$. 
Assume that $W\cap Y_l$ is a proper intersection for every $l\geq e+1$.
Then 
$$\dim f^{-1}(W)=\dim W+e.$$
Moreover, for every irreducible component $Z$ of $f^{-1}(W)$ with $\dim Z=\dim f^{-1}(W)$, we have $f(Z)=W.$
\end{lem}

\proof
Write $W=\sqcup_{e\leq l\leq \dim X} W\cap(Y_l\setminus Y_{l+1}).$
For every $l\geq e+1$, if $W\cap(Y_l\setminus Y_{l+1})\neq \emptyset,$
$$\dim W\cap(Y_l\setminus Y_{l+1})\leq \dim W\cap Y_l=\dim W+\dim Y_l-\dim Y$$
$$\leq \dim W+\dim X-l-1-\dim Y=\dim W+e-l-1.$$
So for $l\geq e+1$, if $W\cap(Y_l\setminus Y_{l+1})\neq \emptyset,$ 
$$\dim f^{-1}(W\cap (Y_l\setminus Y_{l+1}))\leq \dim W+e-1.$$


\medskip

Because $W\cap Y_{e+1}$ is a proper intersection, we have $W\setminus Y_{e+1}\neq \emptyset.$
Then $$\dim f^{-1}(W\setminus Y_{e+1})=\dim W+e.$$
So we have $$\dim f^{-1}(W)=\dim W+e.$$
Let $Z$ be an irreducible component of $f^{-1}(W)$ of $\dim Z=\dim f^{-1}(W)=\dim W+e$.
Then $\dim f(Z)\geq \dim Z-e= \dim W.$ So $f(Z)=W.$
%\medskip
%Now assume that $W$ is a local complete intersection in $Y$.
%Then $f^{-1}(W)$ is a local complete intersection in $X$. 
% So every irreducible component of $f^{-1}(W)$ has dimension at least $e+\dim W.$ 
%So $f^{-1}(W)$ is of pure dimension $\dim W+e.$
\endproof



\proof[Proof of Proposition \ref{prostrleqtot}]
For every $i=1,\dots,n$, since $f(Z_i)=V$ and $Z_i$ is irreducible, 
$\dim Z_i=\dim (Z_i\setminus f^{-1}(Y_{e+1})).$
Since $f^{-1}(V\setminus Y_{e+1})$ is of pure dimension $e+\dim V$,  $\dim Z_i=e+\dim V.$ 

Set $r:=\dim Y-\dim V.$ 
By Proposition \ref{probertini},  there are effective and very ample divisors $H_1,\dots, H_r$ on $Y$ such that 
$H_1\cap \dots \cap H_r$ is a proper intersection and 
$$H_1\cdots H_r=V+W$$
where $W$ is an effective $(\dim Y-r)$-cycle such that 
$V\not\subseteq \Supp \,W$, and such that $W\cap Y_l$ is a proper intersection
for $l\geq e+1$.
We have $$g^*[V]+g^*[W]=[g^*H_1\cdots g^*H_r].$$

By Lemma \ref{lempullbackproper}, $\dim f^{-1}W=\dim V+e.$  
Then $f^{-1}W=g^{-1}W \cap X$ is a proper intersection. Since $B$ is smooth, $f^{-1}W$ is equidimensional. 
By Lemma \ref{lempullbackproper} again, for every irreducible component $R$ of $f^{-1}W$, the image $f(R)$ is an irreducible component of $W.$

We claim that $[f^{-1}W]=g^*[W]\cdot [X]$ as algebraic cycles over $B$. 
In fact, it suffices to prove the restriction of the equality to $U=B\setminus g^{-1}(V)$. Over $Y\setminus V$, $W$ is the complete intersection $H_1\cap \dots \cap H_r$. 
Then the result follows by the fact that the intersection multiplicities are given by length of the local rings. 
This can be obtained by the vanishing of the higher tor-functors in Serre's intersection formula, or as an example of \cite[Proposition 7.1 (b)]{Fulton1984}. 

Similarly, we have the Zariski closure
$$[\overline{f^{-1}(V)\cap f^{-1}(Y\setminus Y_{e+1})}]
=\sum_{i=1}^n m_i[Z_i]$$
as algebraic cycles over $B$. 
The sum gives 
$$[\overline{f^{-1}(H_1\cap \dots \cap H_r)\cap f^{-1}(Y\setminus Y_{e+1})}]=\sum_{i=1}^n m_i[Z_i]+g^*[W]\cdot [X].$$



Note that every irreducible component of 
$f^{-1}(H_1\cap \dots \cap H_r)\cap f^{-1}(Y\setminus Y_{e+1})$ has dimension $\dim V+e.$
So we have 
$$[\overline{f^{-1}(H_1\cap \dots \cap H_r)\cap f^{-1}(Y\setminus Y_{e+1})}]\leq (g^*H_1\cdots g^*H_r\cdot X)^{\prop}$$
in $\CH_{\dim V+e}(B)_\Q$. 

Finally, by Proposition \ref{propropintboundhyp}, since $\sO_{B}(f^*H_i)$ for $ i=1,\dots,r$ are generated by global sections, 
we get 
$$(g^*H_1\cdots g^*H_r\cdot X)^{\prop}\leq g^*H_1\cdots g^*H_r\cdot [X].$$
It follows that
$$\sum_{i=1}^n m_i[Z_i]+g^*[W]\cdot [X]\leq g^*H_1\cdots g^*H_r\cdot [X]=g^*[V]\cdot X+g^*[W]\cdot [X].$$
This concludes the proof.
\endproof






%\begin{pro}\label{probertini}Let $X$ be a projective variety over an algebraically closed field $k$.
%Let $H$ be an ample divisor on $X$.
%Let $V$ be an irreducible subvariety of $X$ of codimension $r\geq 1.$
%Let $Z_1,\dots,Z_m$ be irreducible subvarieties of $X$.
%Let $W$ be a subvariety of $X.$
%Assume that for every $i=1,\dots,m$, $Z_i\not\subseteq V\cup W.$
%Then there are $n\geq 1$ and $H_1,\dots, H_r\in |nH|$ such that 
%\begin{points}
%\item $H_1\cap\dots\cap H_r$ is a proper intersection;
%\item $H_1\cdot\dots\cdot H_r=aV+S$ where $S$ is an effective $(\dim X-r)$-cycle such that $a>0,$
%$V\subseteq \Supp \,S$, $S\cap Z_i, i=1,\dots,m$ is a proper intersection, and for every $i=1,\dots,m$, no irreducible component of $S\cap Z_i$ is contained in $W.$
%\end{points}
%\end{pro}
%
%\proof[Proof of Proposition \ref{probertini}]
%Consider the blowup along $Z$
%$$\pi: X':={\rm Bl}_{Z} X\to X.$$
%Denote by $E:=\pi^{-1}(Z)$ the exceptional divisor of $\pi.$ It is an effective Cartier divisor and $\sO_{X'}(-E)=\pi^{-1}I_Z\cdot \sO_{X'}.$
%Because $-E$ is $\pi$-ample, $\pi^*H-E$ is ample on $X'.$
%
%
%
%
%For every $n\geq 1$, the natural morphism $$\pi^*I_Z\twoheadrightarrow \pi^{-1}I_Z\cdot \sO_{X'}=\sO_{X'}(-E)$$ induces a morphism
%$\tau_n: \sI_V^n\to  \pi_*O(-nE).$
%\begin{lem}\label{lemdirectimageblowup}
%For $n>> 0$, $\tau_n$ is an isomorphism. Moreover, if $X$ is normal, $\tau_1$ is an isomorphism.
%\end{lem}
%We fix $n\geq 1$, such that $\tau_n$ is an isomorphism and $n(\pi^*H-E)$ is very ample on $X'.$
%Then we have an isomorphism $$\id\otimes \tau_n^{-1}: \pi_*\sO_{X'}(n(\pi^*H-E))=\sO_X(nH)\otimes_{\sO_X}\pi_*(\sO_{X'}(-nE))\to \sO_X(nH)\otimes_{\sO_X} I_V^n.$$
%Composition $\id\otimes \tau_n^{-1}$ with the injection $\sO_X(nH)\otimes_{\sO_X} I_V^n\hookrightarrow \sO_X(nH),$
%we get a morphism $$\phi:\pi_*\sO_{X'}(n(\pi^*H-E))\to \sO_X(nH).$$
%
%
%
%
%
%\medskip
%
%For $i=1,\dots,m$, set $Z_i':=\pi^{-1}_\#Z_i$ the strict transform of $Z_i.$ Set $W':=\pi^{-1}W.$
%\begin{lem}\label{lemgensec}There are $D_1,\dots, D_r\in |n(\pi^*H-E)|$
%such that for every $i=1,\dots,m$, no irreducible component of $D_i$ is contained in $E,$
%$D_1\cap \dots \cap D_r$ is a proper intersection and for every irreducible component $R$ of $D_1\cap \dots \cap D_r$, and we have 
%\begin{points}
%\item $R\cap Z_i'$ is a proper intersection for every $i=1,\dots,m$;
%\item no irreducible component of $R\cap Z_i'$ is contained in $E\cup W'.$
%\end{points}
%\end{lem}
%Pick $D_1,\dots, D_r\in |n(\pi^*H-E)|$ as in Lemma \ref{lemgensec}.
%They are zero locus of section $u_1,\dots, u_r\in \sO_{X'}(n(\pi^*H-E)).$
%
%We describe the morphism $$\phi_X: H^0(\sO_{X'}(n(\pi^*H-E)))=H^0(\pi_*\sO_{X'}(n(\pi^*H-E)))\to H^0(\sO_X(nH))$$ induced by $\phi$ as follows.
%Every section $u\in H^0(\sO_{X'}(n(\pi^*H-E)))$ can be viewed as a rational function on $X'$ with 
%$\Div_{X'}(u)+n(\pi^*H-E)\geq 0.$ Then $\phi_X(u)$ is the rational function $(\pi^*)^{-1}(u)\in k(X).$ The divisor $H_u$ on $X$ defined by $\phi_X(u)$ is 
%$\Div_X((\pi^*)^{-1}(u))+nH$, which is effective. 
%Because $(\pi^*)^{-1}(u)$ is contained in the subspace $H^0(\sO_X(nH)\otimes I^n_V)\subseteq H^0(\sO_X(nH))$, 
%$V\subseteq H_u.$
%
%
%Set 
%$s_i:=\phi_X(u_i)\in H^0(\sO_X(nH)), i=1,\dots,r.$  
%The divisor defined by $s_i$ is 
%$$H_i=\Div_X((\pi^*)^{-1}(u))+nH=\pi_*(\Div_X'(u_i))+nH$$
%$$=\pi_*(D_i-n(\pi^*H-E))+nH=\pi_*(D_i+E).$$
%Then $$\Supp\, H_i=\Supp\, \pi_*(D_i+E)\subseteq \pi(D_i)\cup \pi(E)=\pi(D_i)\cup V.$$
%Because no irreducible component of $D_i$ is contained in $E,$ $$\pi(D)\subseteq \Supp\, \pi_*(D_i+E)=\Supp\, H_i.$$
%The discussion in the previous paragraph shows that $V\subseteq H_i.$
%So we have $$\Supp\, H_i=\pi(D_i)\cup V.$$
%Note that, when $r=1$, $V$ may not be contained in $\pi(D_i).$
%
%Note that every irreducible component of $\cap _{i=1}^rH_i$ has codimension at most $r=\codim V.$
%To proof that $\cap _{i=1}^rH_i$ is a complete intersection, we only need to show that $\Supp\, \cap _{i=1}^rH_i$ has pure codimension $r.$
%Because $\pi|_{X'\setminus E}\to X\setminus V$ is an isomorphism and $\pi(E)=V$, 
%we have $$\Supp\, \cap _{i=1}^rH_i= \cap _{i=1}^r\Supp\, H_i=\cap _{i=1}^r(\pi(D_i)\cup V)=V\cup (\cap_{i=1}^m\pi(D_i)\setminus V)$$
%$$=V\cup (\cap_{i=1}^m\pi(D_i\setminus E))=V\cup \pi(\cap_{i=1}^m(D_i\setminus E))=V\cup \pi((\cap_{i=1}^mD_i)).$$
%Because $\cap_{i=1}^mD_i$ is a proper intersection, 
%$V\cup \pi((\cap_{i=1}^mD_i))$ is of pure codimension $r$, hence is a proper intersection.
%Moreover, since no irreducible component of $\cap_{i=1}^mD_i$ is contained in $E$, every irreducible component $Y$ of $\cap_{i=1}^rH_i$ except $V$ takes form 
%$\pi(Y')$ where $Y'$ is an irreducible component of  $\cap_{i=1}^mD_i.$ 
%
%We have $$\pi(Y'\cap Z_i')\subseteq \pi(Y')\cap Z_i$$
%
%
%
%\endproof
%
%\proof[Proof of Lemma \ref{lemdirectimageblowup}]
%The definition of the blowup $X'={\rm Proj}\oplus_{i\geq 0}I^i_V$ where the polarization is given by $\sO_{X'}(-E).$
%So we have $\tau_n:I_V^n\to  \pi_*\sO_{X'}(-nE)$ is injective and  for $n>>0$, $\tau_n$ is an isomorphism. 
%
%Now assume that $X$ is normal.
%By Zariski's main theorem, we may identify $\pi_*\sO_{X'}(-E)$ as an ideal sheaf of $\sO_X$ via  $\pi_*\sO_{X'}(-E)\hookrightarrow \pi_*\sO_{X'}=\sO_X.$
%Every local section of $s$ of $(\pi_*\sO_{X'}(-E))$ is a rational function on $X$ such that $\pi^*s$ vanish on $E.$ So $s$ vanishes on $V.$ Because $V$ is reduced, $s\in I_V.$ It follows that $\tau_1$ is surjective. So $\tau_1$ is an isomorphism.
%\endproof
%
%\proof[Proof of Lemma \ref{lemgensec}]
%\endproof
